\uuid{lVg5}
\titre{Écriture algébrique}
\auteur{bodin ; hanani ; mzari}
\organisation{exo7}
\licence{CC BY-NC-SA 4.0}
\datecreate{2020-01-20}
\niveau{L1}
\module{Algèbre}
\chapitre{Nombres complexes}
\sousChapitre{Forme cartésienne, forme polaire}
\difficulte{1}
\theme{Écritures algébrique et géométrique, écriture algébrique}
\exo7qcmid{lille-1.nombres-complexes.003}

\contenu{
\begin{question}
\qclassification{104.01}
\qtags{motcle=écriture algébrique}
Soit $z$ un nombre complexe de module $2$ et d'argument $\frac{\pi}{4}$. L'écriture algébrique de $z$ est :
\begin{answers}
    \bad{$z= \sqrt 2-i\sqrt 2$}

    \good{$z= \sqrt 2+i\sqrt 2$}

    \bad{$z= 2+2i$}

    \bad{$z= 2-2i$}
\end{answers}
\begin{explanations}
$z=2(\cos\frac{\pi}{4}+i\sin\frac{\pi}{4}) =\sqrt 2+i\sqrt 2 $.
\end{explanations}
\end{question}
}
